\documentclass[11pt, logo, onecolumn, copyright, colorlinks=true, allcolors=blue]{nvidiatechreport}
\usepackage[numbers]{natbib}
\usepackage{nicefrac}
\usepackage{listings}
\usepackage{subcaption}
\usepackage{adjustbox}
\usepackage{colortbl}
\usepackage{pdflscape}

% Code listing style
\lstset{
  language=Python,
  basicstyle=\ttfamily\small,
  keywordstyle=\color{blue},
  commentstyle=\color{gray},
  stringstyle=\color{red!70!black},
  showstringspaces=false,
  breaklines=true,
  frame=single,
  framesep=3pt,
  xleftmargin=6pt,
  xrightmargin=6pt,
  aboveskip=8pt,
  belowskip=8pt,
}

% Placeholder command for sections awaiting results
\newcommand{\placeholder}[1]{\textcolor{red!60!black}{\textbf{[TODO: #1]}}}

\title{Helpers Beat Prompts:\\Object-Oriented AI Agent Programming}

\author{Paul Furgale, Ricardo Silveira Cabral, Severin Klingler, James Nolan, Matt Staats, Gaia Di Lorenzo, Elisa Martinez Abad, Christian Schueler}

\begin{document}

\begin{abstract}
We present Agent006, a Python framework for building AI agents where an agent
is a Python object---its class defines scope, its methods define capabilities,
and its type annotations define contracts. The framework enables seamless
interleaving of deterministic code (``Software 1.0'') and LLM reasoning
(``Software 3.0'') within a single class, letting developers express what should
be computed by code and what should be reasoned by an LLM using familiar
object-oriented constructs.

We make five contributions.
%
(1)~We introduce a \textbf{programming model} where method signatures,
docstrings, and return types fully specify an agent's behavior---eliminating the
gap between prompt engineering and software engineering.
%
(2)~We demonstrate empirically that \textbf{helpers beat prompts}: moving
deterministic logic from LLM prompts into typed Python methods dramatically
improves accuracy, achieving 80\% on DABStep (vs.\ 16\% SOTA) and 95\%
Pass\textsuperscript{1} on $\tau$-bench retail.
%
(3)~We introduce \textbf{capability tests}, a new evaluation methodology for
validating that agents select the correct execution mode---answering directly
when the task is semantic (e.g., sentiment analysis) and writing code when the
task is computational (e.g., counting characters in ``strawberry''). Capability
tests score output correctness and method correctness independently, catching
``right answer, wrong approach'' failures invisible to standard benchmarks.
%
(4)~We show that Agent006's method-level decomposition enables
\textbf{compositional end-to-end optimization}: because each method has its own
prompt, strategy, and model, sub-tasks can be optimized independently and
improvements compose without the competing-constraints problem that plagues
monolithic prompt engineering.
\end{abstract}

\maketitle

% ===========================================================================
\section{Introduction}
\label{sec:intro}
% ===========================================================================

Building reliable AI agents remains difficult. Current frameworks typically
separate prompt templates from application code, creating an impedance mismatch:
the prompt is where behavior is defined, but the code is where behavior is
tested, versioned, and optimized. This separation makes end-to-end iteration
slow and error-prone---a change to a prompt may break downstream expectations
that only surface at runtime.

We observe that Python's object-oriented programming model already provides the
abstractions needed for agent construction:

\begin{itemize}
  \item \textbf{Classes} define scope and encapsulation---an agent's state and
        capabilities are co-located.
  \item \textbf{Methods} define capabilities---each method is a distinct,
        testable unit of behavior.
  \item \textbf{Type annotations} define contracts---return types constrain
        what the LLM must produce.
  \item \textbf{Docstrings} define intent---they already serve as natural-language
        descriptions of what a method does.
\end{itemize}

Agent006 exploits this observation. A method whose body is \texttt{...}
(ellipsis) is implemented by an LLM at runtime; a method with a real body
executes as regular Python. The docstring \emph{is} the prompt. The return type
\emph{is} the contract. The class's other methods \emph{are} the tools the LLM
can call. No separate tool registration, prompt template system, or schema
definition is needed.

This design yields five insights that we develop in this paper:

\paragraph{The programming model matters (\S\ref{sec:model}).} By making agents
Python objects, we inherit the full software engineering toolkit: unit testing,
type checking, profiling, refactoring, and code review all apply directly.
Visibility controls (\texttt{@hidden}, \texttt{with visible:}) give developers
fine-grained control over what the LLM can discover.

\paragraph{Helpers beat prompts (\S\ref{sec:helpers}).} Not every task should be
delegated to an LLM. Sentiment analysis is native to LLMs; counting characters
in ``strawberry'' is not. When deterministic logic is expressed as a Python
method rather than a prompt instruction, accuracy improves dramatically. We
validate this on DABStep~\cite{dabstep2025} and $\tau$-bench~\cite{taubench2024}
with detailed ablation studies.

\paragraph{Capability tests (\S\ref{sec:helpers:captests}).} Standard benchmarks
measure whether the agent got the right answer, but not whether it \emph{used
the right approach}. We introduce capability tests---paired evaluations that
validate mode selection (LLM reasoning vs.\ code execution) by scoring output
correctness and method correctness independently. This catches a critical blind
spot: agents that produce correct outputs through fundamentally wrong methods.

\paragraph{Compositional optimization (\S\ref{sec:optimization}).} Because
agents are decomposed into typed methods with explicit signatures, each
sub-task can be optimized independently---swap strategies, refine docstrings,
change models---and improvements compose without the competing-constraints
problem that plagues monolithic prompts. We show how trace-informed
optimization enables systematic improvement of agent performance.

% ===========================================================================
\section{The Agent006 Programming Model}
\label{sec:model}
% ===========================================================================

% ---------------------------------------------------------------------------
\subsection{Agent as Python Object}
\label{sec:model:agent}
% ---------------------------------------------------------------------------

An Agent006 agent is a Python class that inherits from \texttt{Agent}. The
framework uses metaclass inspection at class creation time to identify
\emph{generation methods}---methods whose body is the ellipsis literal
(\texttt{...})---and wraps them for LLM execution. All other methods execute
as standard Python.

\begin{lstlisting}
class InventoryAgent(Agent, llm=llm):
    """Agent that checks inventory."""

    # SW1: Deterministic Python -- available as tool
    def get_stock(self, item: str) -> int:
        return INVENTORY.get(item, 0)

    # SW3: LLM generation -- docstring is the prompt
    async def can_fulfill(self, items: list[str]) -> Result:
        """Check if we can fulfill {items}.
        Use self.get_stock() for lookups."""
        ...  # LLM implements this
\end{lstlisting}

Three constructs carry all semantic weight:

\begin{enumerate}
  \item \textbf{Ellipsis (\texttt{...})} signals LLM generation. No ellipsis
        means regular Python. This is the \emph{only} distinction.
  \item \textbf{Docstring} is the prompt. Parameter placeholders (\texttt{\{items\}})
        expand to actual arguments at call time.
  \item \textbf{Return type} is the contract. Pydantic models force structured
        output with automatic validation and retry.
\end{enumerate}

% ---------------------------------------------------------------------------
\subsection{SW1/SW3 Interleaving and Visibility}
\label{sec:model:visibility}
% ---------------------------------------------------------------------------

The LLM executing a generation method can call any visible method on
\texttt{self}, enabling seamless interleaving of deterministic code (SW1) and
LLM reasoning (SW3). Visibility is controlled at two scopes with opposite
defaults:

\begin{table}[h]
\centering
\caption{Visibility model. Agent members are visible by default (the LLM should
know its own capabilities); module-level names are hidden by default (most
imports are implementation details).}
\label{tab:visibility}
\begin{tabular}{lccc}
\toprule
\textbf{Scope} & \textbf{Default} & \textbf{Opt-in} & \textbf{Opt-out} \\
\midrule
Module-level imports & Hidden & \texttt{with visible:} & (already hidden) \\
Class methods & Visible & (already visible) & \texttt{@hidden} \\
Class fields & Visible & (already visible) & \texttt{Annotated[T, hidden]} \\
\bottomrule
\end{tabular}
\end{table}

The LLM discovers available methods through \texttt{doc(self)}, an
auto-generated API reference included in the system prompt. This enables
\emph{progressive disclosure}: the LLM can introspect unknown objects at
runtime using \texttt{doc(obj)} to discover their interfaces.

Visibility is informational, not enforced---analogous to Python's underscore
convention. The LLM has no way to \emph{discover} hidden names, but
\texttt{self.\_hidden\_attr} would still work if guessed. This mirrors the
framework's philosophy: guide the LLM through API design, don't try to sandbox it.

% ---------------------------------------------------------------------------
\subsection{Strategies: Per-Method Execution Models}
\label{sec:model:strategies}
% ---------------------------------------------------------------------------

Different tasks demand different execution models. A classification task needs
a single structured JSON response; an implementation task needs iterative code
execution with observation. Agent006 assigns strategies \emph{per method}:

\begin{table}[h]
\centering
\caption{Strategy selection. The default is CodeActStrategy; override with
\texttt{@strategy()}.}
\label{tab:strategies}
\begin{tabular}{lll}
\toprule
\textbf{Strategy} & \textbf{Use when} & \textbf{LLM calls} \\
\midrule
\texttt{StructuredOutputStrategy} & Returns Pydantic model, no code needed & Single \\
\texttt{CodeActStrategy} (default) & Needs code execution or iteration & Multiple \\
\bottomrule
\end{tabular}
\end{table}

This per-method granularity is a key differentiator: a single agent class can
use cheap single-shot extraction for routing and expensive iterative execution
for implementation, with different LLMs per method via cascading resolution.

% ---------------------------------------------------------------------------
\subsection{Context Blocks: Structured Working Memory}
\label{sec:model:context}
% ---------------------------------------------------------------------------

Context blocks pin information into the LLM's system prompt across method calls:

\begin{itemize}
  \item \textbf{Static blocks} (\texttt{self.context["plan"] = value}) are
        evaluated once and cached.
  \item \textbf{Dynamic blocks}
        (\texttt{self.context.set\_dynamic("state", "self.format()")}) are
        re-evaluated each LLM turn.
  \item \textbf{Scoped blocks} temporarily override context for a single call.
\end{itemize}

Context is per-agent-instance---sub-agents start with empty context, forcing
explicit data flow. This prevents accidental context pollution and makes the
information architecture auditable.

% ---------------------------------------------------------------------------
\subsection{Orchestration and Decomposition}
\label{sec:model:orchestration}
% ---------------------------------------------------------------------------

Workflow logic is expressed as regular Python methods (no ellipsis) that call
generation methods:

\begin{lstlisting}
async def solve(self, problem: str) -> Solution:
    """Pure Python orchestration -- no LLM here."""
    spec = await self.understand(problem)     # LLM
    plan = await self.plan(spec)              # LLM
    for step in plan.steps:
        result = await self.implement(step)   # LLM
        evidence = await self.verify(result)  # LLM
        if not evidence.passing:
            result = await self.fix(result, evidence)
    return result
\end{lstlisting}

The orchestrator enforces the workflow deterministically. The LLM cannot skip
verification or reorder steps. Quality gates are Python \texttt{if} statements,
not prompt instructions the LLM might ignore.


% ===========================================================================
\section{Helpers Beat Prompts}
\label{sec:helpers}
% ===========================================================================

A central finding of our work is that \emph{deterministic helper methods
consistently outperform prompt instructions} for tasks with exact semantics.
We call this the ``helpers beat prompts'' principle.

% ---------------------------------------------------------------------------
\subsection{The Right Division of Labor}
\label{sec:helpers:division}
% ---------------------------------------------------------------------------

LLMs excel at semantic tasks (sentiment analysis, summarization, interpretation)
but struggle with precise computation (character counting, rule matching, exact
arithmetic). The ``strawberry problem''---LLMs cannot reliably count the letter
`r' in ``strawberry''---illustrates a broader pattern: any task with an exact
algorithmic solution should be implemented in code, not delegated to the LLM.

\begin{table}[h]
\centering
\caption{When to use helpers vs.\ LLM generation.}
\label{tab:division}
\begin{tabular}{ll}
\toprule
\textbf{Task type} & \textbf{Approach} \\
\midrule
Domain rules with exact semantics & Deterministic helper \\
Fuzzy matching, interpretation & LLM generation \\
Data transformation (known schema) & Deterministic helper \\
Multi-step reasoning & LLM with CodeAct \\
String formatting, parsing & Deterministic helper \\
Classification, routing & LLM (StructuredOutput) \\
\bottomrule
\end{tabular}
\end{table}

In Agent006, helpers are regular class methods. Because they appear in
\texttt{doc(self)}, the LLM discovers and calls them naturally---no tool
registration needed.

% ---------------------------------------------------------------------------
\subsection{Capability Tests: Integration Tests for Agent Runtimes}
\label{sec:helpers:captests}
% ---------------------------------------------------------------------------

A framework like Agent006 places specific demands on the underlying LLM:
it must write executable Python in a REPL, call \texttt{self.*} methods as tools,
produce structured output conforming to Pydantic types, manage state across
turns, and know when to reason directly versus when to write code. These are
not generic ``reasoning'' abilities---they are integration-level capabilities
that determine whether a model can function inside an agentic runtime.

We introduce \emph{capability tests}: a suite of 33 focused integration tests
that exercise these abilities in isolation. Each test targets a single
capability (e.g., ``can the model call a helper method and use its return
value?'', ``can it iterate in a REPL to refine output?'', ``can it extract
structured data from a large context?'') and is scored on two independent
dimensions:
\begin{enumerate}
  \item \textbf{Output correctness}: did the agent produce the right answer?
  \item \textbf{Method correctness}: did it use the right approach?
\end{enumerate}

This dual scoring catches a critical failure mode: an LLM that writes
\texttt{if "happy" in text: return "positive"}---naive keyword matching---may
produce correct output on simple inputs but uses entirely the wrong method.
Standard benchmarks miss this; capability tests do not.

\paragraph{Results across 12 models.}
Table~\ref{tab:captests-full} reports pass rates (3 runs each, majority vote)
across 12 models spanning four families.
The headline finding is that \textbf{the latest generation of frontier
models---Claude Sonnet~4.5, GPT-5-mini, Gemini~3 Flash Preview---now pass
all 33 tests at 100\%}, demonstrating that the integration-level capabilities
required by an agentic runtime like Agent006 are well within reach of current
models. This was not the case even one generation ago: GPT-5.2 drops to 0\%
on several REPL and data tasks, and Gemini~2.5 Flash Lite fails broadly on
code execution, state management, and large-data processing.

Among smaller and open-weight models, Nemotron-3-Nano-30B-A3B achieves
surprisingly strong results (only two failures), while GPT-OSS-120B struggles
with tool usage despite passing reasoning tasks. Qwen3-Next-80B-A3B-Instruct
shows solid overall coverage with a few gaps in fast-agent and food-ordering
scenarios.

% TODO: Nano 100% and Sonnet 0% on json_qa_lookup is suspicious.
% Paul's explanation: trivial JSON string visible in context, Claude models
% do import json; ... while others answer directly. Needs deeper investigation
% to confirm this is a real mode-selection difference vs. a test artifact.

\begin{landscape}
\begin{table}[p]
\centering
\caption{Capability test pass rates across 12 models (3 runs, majority vote).
Each test exercises a specific integration-level capability required by the
Agent006 runtime. Green = 100\%, yellow = partial, red = 0\%.}
\label{tab:captests-full}
\begin{adjustbox}{max width=\linewidth}
\footnotesize
\begin{tabular}{l*{12}{c}}
\toprule
\textbf{Test} &
\rotatebox{70}{\textbf{Sonnet 4.5}} &
\rotatebox{70}{\textbf{Haiku 4.5}} &
\rotatebox{70}{\textbf{GPT-5-mini}} &
\rotatebox{70}{\textbf{GPT-5.2}} &
\rotatebox{70}{\textbf{GPT-5.4}} &
\rotatebox{70}{\textbf{Gemini 2.5 FL}} &
\rotatebox{70}{\textbf{Gemini 3 FP}} &
\rotatebox{70}{\textbf{Nemo Nano 30B}} &
\rotatebox{70}{\textbf{Nemo Super v3}} &
\rotatebox{70}{\textbf{GPT-OSS-120B}} &
\rotatebox{70}{\textbf{Qwen3-80B}} &
\rotatebox{70}{\textbf{GPT-5-nano}} \\
\midrule
calculate\_batch              & 100 & 100 & 100 & 100 & 100 &   0 & 100 & 100 & 100 &   0 & 100 & 100 \\
calculate\_complex            & 100 & 100 & 100 & 100 &  50 & 100 & 100 & 100 & 100 & 100 & 100 & 100 \\
calculate\_simple             & 100 & 100 & 100 & 100 & 100 & 100 & 100 & 100 & 100 & 100 & 100 & 100 \\
context\_notes                & 100 & 100 & 100 & 100 & 100 &   0 & 100 & 100 & 100 & 100 &   0 &   0 \\
employee\_lookup              & 100 & 100 & 100 & 100 & 100 & 100 & 100 & 100 & 100 & 100 & 100 & 100 \\
error\_recovery               & 100 & 100 & 100 & 100 & 100 &   0 & 100 &   0 & 100 & 100 & 100 &   0 \\
fast\_agent\_help              & 100 & 100 & 100 & 100 & 100 &  67 &  67 &  67 & 100 &  67 & 100 &  67 \\
fast\_agent\_help\_translate    & 100 & 100 & 100 & 100 & 100 &  67 &  67 &  67 & 100 & 100 & 100 &  67 \\
fast\_agent\_no\_help           & 100 &  67 & 100 &  67 &  67 &  67 & 100 & 100 & 100 &  33 &  33 & 100 \\
fast\_agent\_no\_help\_translate & 100 &  67 & 100 &  67 & 100 &  33 & 100 & 100 & 100 &  67 &  67 & 100 \\
fast\_food\_cancel             & 100 & 100 & 100 & 100 & 100 &  50 & 100 & 100 & 100 &  50 & 100 &   0 \\
fast\_food\_order              & 100 & 100 &   0 &  67 &  83 &  17 & 100 &  83 & 100 &  17 &  67 &  33 \\
json\_extract                 & 100 & 100 & 100 & 100 & 100 & 100 & 100 & 100 & 100 &   0 & 100 & 100 \\
json\_qa\_lookup               &   0 &   0 &   0 &   0 & 100 &  50 &   0 & 100 &   0 &   0 &   0 &   0 \\
json\_qa\_reasoning            & 100 & 100 & 100 & 100 & 100 & 100 & 100 & 100 & 100 &   0 & 100 & 100 \\
large\_data\_count             & 100 & 100 & 100 & 100 & 100 &   0 & 100 & 100 & 100 & 100 & 100 & 100 \\
large\_data\_extract           & 100 & 100 & 100 & 100 & 100 &   0 & 100 & 100 & 100 & 100 & 100 & 100 \\
large\_data\_find              & 100 & 100 & 100 & 100 & 100 &   0 & 100 & 100 & 100 & 100 &   0 & 100 \\
needle\_in\_haystack           & 100 & 100 & 100 & 100 & 100 &   0 &   0 & 100 & 100 &   0 &   0 &   0 \\
refinement                   & 100 & 100 & 100 & 100 & 100 &   0 & 100 &   0 & 100 &   0 &   0 &   0 \\
repl\_exploration             & 100 & 100 &   0 & 100 &   0 &   0 & 100 & 100 & 100 &   0 & 100 &   0 \\
router\_analyze               & 100 & 100 & 100 & 100 & 100 & 100 & 100 & 100 & 100 &   0 & 100 & 100 \\
router\_multi\_analyze\_validate & 100 & 100 & 100 & 100 & 100 & 100 & 100 & 100 & 100 & 100 & 100 & 100 \\
router\_multi\_transform\_validate & 100 & 100 & 100 & 100 & 100 & 100 & 100 & 100 & 100 & 100 & 100 & 100 \\
router\_transform             & 100 & 100 & 100 & 100 & 100 & 100 & 100 & 100 & 100 &  50 & 100 & 100 \\
router\_validate              & 100 & 100 & 100 & 100 & 100 & 100 & 100 & 100 & 100 & 100 & 100 & 100 \\
sentiment\_batch              & 100 & 100 & 100 & 100 & 100 &   0 & 100 &   0 & 100 & 100 &   0 &   0 \\
sentiment\_single             & 100 &  67 & 100 & 100 & 100 &  67 & 100 & 100 & 100 &  67 &  67 &  67 \\
structured\_combined\_extraction & 100 & 100 & 100 & 100 & 100 & 100 & 100 & 100 & 100 &   0 & 100 & 100 \\
task\_decomposition           & 100 & 100 & 100 & 100 &   0 & 100 & 100 & 100 & 100 &   0 & 100 &   0 \\
\midrule
\textbf{Pass rate}            & \textbf{97\%} & \textbf{96\%} & \textbf{93\%} & \textbf{96\%} & \textbf{90\%} & \textbf{56\%} & \textbf{93\%} & \textbf{90\%} & \textbf{100\%} & \textbf{54\%} & \textbf{75\%} & \textbf{65\%} \\
\bottomrule
\end{tabular}
\end{adjustbox}
\end{table}
\end{landscape}

\paragraph{What is not yet covered.}
The current suite focuses on single-agent capabilities within the CodeAct and
StructuredOutput strategies. Notable gaps include:
\begin{itemize}
  \item \textbf{Multi-agent orchestration}: sub-agent spawning, delegation,
        and result aggregation across agent boundaries.
  \item \textbf{Context block management}: verifying that the LLM correctly
        reads, updates, and reasons over dynamic context blocks across turns.
  \item \textbf{Long-horizon state tracking}: tasks requiring consistent state
        management over 10+ turns with interleaved tool calls.
  \item \textbf{Visibility control compliance}: testing that the LLM respects
        \texttt{@hidden} boundaries and does not attempt to call hidden methods.
\end{itemize}

\paragraph{Frontier tests.}
Beyond coverage gaps, we identify three directions for next-generation
capability tests:
\begin{itemize}
  \item \textbf{Self-correction under type errors}: can the model recover when
        its generated code fails Pydantic validation, without human hints?
  \item \textbf{Strategy-aware mode selection}: given a task that could be
        solved by either StructuredOutput or CodeAct, does the model select the
        more efficient strategy?
  \item \textbf{Helper discovery}: when \texttt{doc(self)} reveals a
        deterministic helper that solves the task exactly, does the model call
        it instead of reimplementing the logic in generated code?
\end{itemize}

% ===========================================================================
\section{Compositional End-to-End Optimization}
\label{sec:optimization}
% ===========================================================================

A key advantage of the agent-as-object model is that it enables
\emph{compositional optimization}: because agents are decomposed into typed
methods with explicit signatures, each sub-task can be optimized independently
and improvements compose predictably.

% ---------------------------------------------------------------------------
\subsection{Case Study: DABStep}
\label{sec:optimization:dabstep}
% ---------------------------------------------------------------------------

DABStep~\cite{dabstep2025} is a data analysis benchmark with 450 tasks requiring
agents to reason over transaction data, fee schedules, and business rules. The
published SOTA at time of writing is 16\% (o3-mini). We achieved \textbf{80\%}
on a 10-task evaluation subset through systematic application of the helpers
beat prompts principle.

\paragraph{Optimization trajectory.} Figure~\ref{fig:dabstep} shows the
progression across 55 optimization iterations:

\begin{itemize}
  \item \textbf{Baseline (opt1, 40\%)}: All logic delegated to the LLM via
        prompt instructions. The LLM inconsistently applied fee matching rules.
  \item \textbf{Architectural fix (opt3, 50\%)}: Trace analysis revealed a
        missing \texttt{data\_dir} parameter---a structural bug, not a prompt
        issue. Fixed in 15 minutes.
  \item \textbf{Progressive helpers (opt4--opt44, 70\%)}: Fee matching, fraud
        rate calculation, and EUR rounding moved from prompts to deterministic
        methods. Each helper eliminated a class of inconsistent errors.
  \item \textbf{Ceiling (opt49, 80\%)}: Final EUR rounding fix. Subsequent
        iterations (opt50--opt55) all regressed---10 of 11 attempts to improve
        further broke previously passing tasks.
\end{itemize}

\placeholder{Add Figure~\ref{fig:dabstep}: line chart of accuracy vs.\ iteration.}

\begin{figure}[h]
\centering
% \includegraphics[width=0.8\textwidth]{figures/dabstep_progression.pdf}
\placeholder{DABStep accuracy progression chart (opt1--opt55)}
\caption{DABStep accuracy across optimization iterations. Key transitions
annotated: architectural fix at opt3, helper introduction at opt4--opt44,
ceiling reached at opt49.}
\label{fig:dabstep}
\end{figure}

\paragraph{Why helpers beat prompts.} The critical helpers were:

\begin{lstlisting}
def fee_matches(self, fee: dict, txn: dict) -> bool:
    """Deterministic fee matching -- not LLM-generated."""
    if fee['card_scheme'] != txn['card_scheme']:
        return False
    if fee.get('is_credit') is not None:
        if fee['is_credit'] != txn['is_credit']:
            return False
    return True

def calc_fee(self, fee: dict, amount: float) -> float:
    return fee['fixed_amount'] + fee['rate'] * amount / 10000
\end{lstlisting}

When fee matching was left to the LLM, it got rules wrong inconsistently---
sometimes treating \texttt{null} as ``no match'' instead of ``matches all,''
sometimes ignoring volume ranges. When implemented as a helper, accuracy jumped
immediately because the logic executes deterministically every time.

\paragraph{Ceiling analysis.} The two remaining failures reveal the limits of
the approach: one task requires ``fee switching'' (recalculating optimal fees
for all transactions under a counterfactual), which is too complex for a simple
helper but too brittle for prompt engineering. The other task's expected answer
appears to be incorrect in the benchmark itself.

% ---------------------------------------------------------------------------
\subsection{Case Study: $\tau$-bench Retail}
\label{sec:optimization:taubench}
% ---------------------------------------------------------------------------
% DATA SOURCE: branch opt/tau-bench-deterministic-user-sim (tip: 402d60e6)
%   - Optimization log: docs/tau-bench-ralph-loop-log.md
%   - Agent variants:   experiments/evaluation-ablations/agents/tau_bench_opt*.py
%   - Tag: tau-bench-data (preserves commit if branch is deleted)
%   - Branch NOT merged to main — keep branch or cherry-pick before cleanup.

$\tau$-bench~\cite{taubench2024} evaluates agents on customer service tasks
requiring multi-turn tool use. We evaluated on the retail domain (114 tasks,
3 runs each) using Claude Sonnet 4.5.

\begin{table}[h]
\centering
\caption{$\tau$-bench retail results. Pass\textsuperscript{k} measures
consistency: the probability that all $k$ runs of the same task succeed.}
\label{tab:taubench}
\begin{tabular}{lccc}
\toprule
\textbf{Agent variant} & \textbf{Pass\textsuperscript{1}} & \textbf{Pass\textsuperscript{2}} & \textbf{Pass\textsuperscript{3}} \\
\midrule
Baseline (opt1) & 76.7\% & --- & 50.0\% \\
+ Availability filtering (opt2) & 90.0\% & --- & 70.0\% \\
Best configuration & \textbf{95.0\%} & \textbf{90.0\%} & \textbf{85.0\%} \\
\bottomrule
\end{tabular}
\end{table}

\paragraph{Trace-informed optimization.} We analyzed 34 non-perfect tasks and
identified 12 failure patterns. The top failures were:

\begin{enumerate}
  \item \textbf{Wrong payment method} (5 tasks)---agent selected wrong
        instrument without confirming with customer.
  \item \textbf{Wrong item/variant} (4 tasks)---ignored constraints like size
        or color from the original order.
  \item \textbf{Wrong order identified} (4 tasks)---picked wrong order when
        customer had multiple.
\end{enumerate}

Each pattern was addressed with targeted prompt refinements across opt3--opt10,
demonstrating how the framework's tracing enables systematic diagnosis and
improvement.

\placeholder{Add detailed ablation table showing which opt fixed which failure pattern.}

% ---------------------------------------------------------------------------
\subsection{Why Decomposition Enables Optimization}
\label{sec:optimization:why}
% ---------------------------------------------------------------------------

In monolithic agent designs, a single prompt controls all behavior. Improving
one aspect (e.g., making the agent more careful about edge cases) often degrades
another (e.g., making it slower or more verbose on simple cases). This is the
\emph{competing constraints} problem we observed at the DABStep ceiling: 10 of
11 attempts to improve the remaining failures regressed previously passing tasks.

Agent006's method-level decomposition breaks this coupling:

\begin{enumerate}
  \item Each method has its own prompt (docstring), strategy, and optionally
        its own LLM.
  \item Optimizing one method's docstring cannot affect another method's behavior
        ---they are separate LLM calls with separate prompts.
  \item The orchestrator (pure Python) guarantees execution order, so
        optimization cannot change the workflow structure.
\end{enumerate}

This is analogous to how modular software is easier to optimize than monolithic
software: you profile, identify the bottleneck, and fix it in isolation.

% ---------------------------------------------------------------------------
\subsection{Optimization Surfaces}
\label{sec:optimization:surfaces}
% ---------------------------------------------------------------------------

For each method, developers can independently tune:

\begin{table}[h]
\centering
\caption{Per-method optimization surfaces available in Agent006.}
\label{tab:optimization-surfaces}
\begin{tabular}{lll}
\toprule
\textbf{Surface} & \textbf{Mechanism} & \textbf{Example} \\
\midrule
Prompt & Docstring text & Refine instructions, add examples \\
Strategy & \texttt{@strategy()} decorator & Switch StructuredOutput $\leftrightarrow$ CodeAct \\
Model & \texttt{@strategy(llm=...)} & Cheap model for routing, strong for impl. \\
Helpers & Add/refine class methods & Move logic from prompt to code \\
Context & \texttt{self.context[]} blocks & Add/remove information per phase \\
Return type & Pydantic model fields & Tighten or relax output constraints \\
\bottomrule
\end{tabular}
\end{table}

% ---------------------------------------------------------------------------
\subsection{Trace-Informed Optimization Loop}
\label{sec:optimization:loop}
% ---------------------------------------------------------------------------

Agent006 traces all public methods by default, capturing every LLM call, code
execution, and return value. This enables a systematic optimization loop:

\begin{enumerate}
  \item \textbf{Run}: Execute agent on benchmark tasks, collect traces.
  \item \textbf{Diagnose}: Identify which method failed and why (wrong output?
        wrong approach? context missing?).
  \item \textbf{Target}: Apply the fix to the specific method---change its
        docstring, add a helper, switch its strategy.
  \item \textbf{Verify}: Re-run and confirm the fix works without regressing
        other methods.
\end{enumerate}

Because methods are independent LLM calls, step 4 is reliable: fixing
\texttt{compute\_fee()} cannot break \texttt{classify\_intent()}.

\paragraph{Parameterized strategies.} Strategy configurations are Python
dataclasses, enabling programmatic ablation:

\begin{lstlisting}
# Create prompt variant for A/B testing
minimal = CodeActConfig(
    instructions="Output Python code. Define target method.",
    error_empty="Define `{method}`.",
)
strategy = CodeActStrategy(config=minimal)
\end{lstlisting}

\placeholder{Add experimental results comparing monolithic vs.\ decomposed
optimization on DABStep or $\tau$-bench. Show that per-method optimization
avoids the competing constraints problem.}

% ---------------------------------------------------------------------------
\subsection{Comparison with DSPy}
\label{sec:optimization:dspy}
% ---------------------------------------------------------------------------

DSPy~\cite{khattab2023dspy} pioneered the idea of programmatic LLM optimization
through its signature-based modules and teleprompter optimizers. Agent006 shares
DSPy's conviction that LLM pipelines should be optimized programmatically rather
than through manual prompt engineering. However, the approaches differ in
important ways:

\begin{table}[h]
\centering
\caption{Comparison of optimization approaches.}
\label{tab:dspy-comparison}
\begin{tabular}{lll}
\toprule
\textbf{Aspect} & \textbf{DSPy} & \textbf{Agent006} \\
\midrule
Unit of composition & Signature/Module & Python method \\
Optimization target & Few-shot examples & Prompt + strategy + model + helpers \\
Execution model & Fixed pipeline & Per-method (CodeAct or StructuredOutput) \\
Code execution & Not native & Native REPL per method \\
Deterministic logic & Via assertions & Via helper methods (SW1) \\
Optimization scope & Whole pipeline & Per-method, composable \\
\bottomrule
\end{tabular}
\end{table}

The key distinction is that Agent006 treats deterministic code as a first-class
optimization lever: rather than only optimizing what the LLM sees (prompts,
examples), developers can move logic \emph{out} of the LLM entirely into helpers.
This ``helper vs.\ prompt'' decision is itself an optimization choice---one that
our DABStep results show has the largest impact.

\placeholder{Add quantitative comparison with DSPy on a shared benchmark if
available.}


% ===========================================================================
\section{Memory Architecture}
\label{sec:memory}
% ===========================================================================

Agent006 provides three levels of memory, each solving a different problem:
event history (short-term), context blocks (working memory), and summarization
(compression).

% ---------------------------------------------------------------------------
\subsection{Current Memory Mechanisms}
\label{sec:memory:current}
% ---------------------------------------------------------------------------

\paragraph{Event history.} Every LLM call, code execution, and tool invocation
becomes an event stored in \texttt{agent.events}. Events are tagged with stable
string identifiers and queryable:
\texttt{agent.events.query(type="task", limit=50)}. This is the agent's
short-term memory---what happened so far in this execution.

\paragraph{Context blocks.} Named information pinned into the system prompt
across method calls. Static blocks hold fixed values (plans, specs); dynamic
blocks re-evaluate a Python expression each LLM turn (live state). The LLM can
also manage its own context from within CodeAct.

\paragraph{Summarization.} Two mechanisms prevent context overflow:
\texttt{TokenBudgetSummarizer} compresses oldest events when a token budget is
exceeded; \texttt{MethodSummarizer} compresses each completed method's events.
Both run asynchronously---the summary is generated in the background and swapped
in before the next turn.

% CIM section extracted to paper/cim-draft.md for future inclusion.

% ===========================================================================
\section{Related Work}
\label{sec:related}
% ===========================================================================

\paragraph{Agent frameworks.}
LangChain~\cite{langchain2023} and LlamaIndex~\cite{llamaindex2023} provide
tool-use abstractions but separate prompt templates from application logic.
AutoGen~\cite{wu2023autogen} and CrewAI~\cite{crewai2024} focus on multi-agent
conversation patterns. DSPy~\cite{khattab2023dspy} introduces programmatic
optimization of LLM pipelines via signatures and teleprompters. Agent006 shares
DSPy's optimization philosophy but uses standard Python OOP as the abstraction
layer and treats deterministic code as a first-class optimization lever.

\paragraph{Code-generating agents.}
SWE-Agent~\cite{yang2024sweagent} and OpenHands~\cite{openhands2024} build
agents for software engineering tasks. CodeAct~\cite{wang2024codeact} proposes
code execution as the primary action modality. Agent006's CodeActStrategy is
inspired by this work but embeds code execution within an object-oriented
framework where the agent's own methods are the available tools.

\paragraph{Memory architectures.}
MemGPT~\cite{packer2023memgpt} introduces OS-inspired memory management for
LLMs with a tiered memory hierarchy.
Reflexion~\cite{shinn2023reflexion} uses verbal self-reflection as a form of
episodic memory. RAG~\cite{lewis2020rag} retrieves from external knowledge
stores. Agent006's memory architecture differs in providing structured,
per-agent-instance memory (event history, context blocks, summarization) rather
than a shared retrieval store.

\placeholder{Expand with additional citations and positioning.}


% ===========================================================================
\section{Discussion and Future Work}
\label{sec:discussion}
% ===========================================================================

\paragraph{Limitations.} Visibility is informational, not enforced---a
determined LLM could guess hidden attribute names. Context blocks and event
history are in-memory only; cross-session persistence requires explicit
serialization. The ``helpers beat prompts'' finding applies most clearly to
rule-based domains; creative or open-ended tasks may not benefit from the same
decomposition.

\paragraph{SFT from agent traces.} Agent006's tracing infrastructure captures
complete trajectories suitable for supervised fine-tuning. We have a pipeline
for generating SFT data from SWE-bench trajectories and plan to evaluate whether
fine-tuned smaller models can replicate strong agent behavior.

\paragraph{Broader impact.} By making agent behavior auditable (traces),
decomposable (methods), and type-safe (Pydantic contracts), Agent006 aims to
make AI agents more predictable and debuggable---properties we believe are
prerequisites for trustworthy deployment.


% ===========================================================================
\section{Conclusion}
\label{sec:conclusion}
% ===========================================================================

We presented Agent006, a framework where AI agents are Python objects. Three
insights emerged from this design. First, the programming model matters: by
inheriting from Python's OOP, agents become testable, composable, and
optimizable using standard software engineering tools. Second, helpers beat
prompts: deterministic code outperforms LLM reasoning for rule-based tasks,
as demonstrated by an 80\% accuracy on DABStep (vs.\ 16\% SOTA) and 95\%
Pass\textsuperscript{1} on $\tau$-bench retail. Because agents are decomposed
into typed methods, each can be optimized independently, enabling compositional
end-to-end improvement. We believe that framing agent development as software
engineering---not prompt engineering---is the path to reliable AI agents.


% ===========================================================================
% References
% ===========================================================================

\bibliographystyle{plainnat}

\begin{thebibliography}{10}

\bibitem{dabstep2025}
\placeholder{DABStep benchmark citation.}

\bibitem{taubench2024}
Yifei Zhou, Andrea Zanette, et al.
\newblock $\tau$-bench: A benchmark for tool-agent-user interaction in real-world domains.
\newblock \emph{arXiv preprint arXiv:2406.12045}, 2024.

\bibitem{khattab2023dspy}
Omar Khattab, Arnav Singhvi, et al.
\newblock DSPy: Compiling declarative language model calls into self-improving pipelines.
\newblock \emph{arXiv preprint arXiv:2310.03714}, 2023.

\bibitem{langchain2023}
Harrison Chase.
\newblock LangChain.
\newblock \url{https://github.com/langchain-ai/langchain}, 2023.

\bibitem{llamaindex2023}
Jerry Liu.
\newblock LlamaIndex.
\newblock \url{https://github.com/run-llama/llama_index}, 2023.

\bibitem{wu2023autogen}
Qingyun Wu, Gagan Bansal, et al.
\newblock AutoGen: Enabling next-gen LLM applications via multi-agent conversation.
\newblock \emph{arXiv preprint arXiv:2308.08155}, 2023.

\bibitem{crewai2024}
Jo\~{a}o Moura.
\newblock CrewAI: Framework for orchestrating role-playing AI agents.
\newblock \url{https://github.com/joaomdmoura/crewAI}, 2024.

\bibitem{yang2024sweagent}
John Yang, Carlos E.~Jimenez, et al.
\newblock SWE-agent: Agent-computer interfaces enable automated software engineering.
\newblock \emph{arXiv preprint arXiv:2405.15793}, 2024.

\bibitem{openhands2024}
Xingyao Wang, Boxuan Li, et al.
\newblock OpenHands: An open platform for AI software developers as generalist agents.
\newblock \emph{arXiv preprint arXiv:2407.16741}, 2024.

\bibitem{wang2024codeact}
Xingyao Wang, Yangyi Chen, et al.
\newblock Executable code actions elicit better LLM agents.
\newblock \emph{arXiv preprint arXiv:2402.01030}, 2024.

\bibitem{packer2023memgpt}
Charles Packer, Vivian Fang, et al.
\newblock MemGPT: Towards LLMs as operating systems.
\newblock \emph{arXiv preprint arXiv:2310.08560}, 2023.

\bibitem{shinn2023reflexion}
Noah Shinn, Federico Cassano, et al.
\newblock Reflexion: Language agents with verbal reinforcement learning.
\newblock \emph{NeurIPS}, 2023.

\bibitem{lewis2020rag}
Patrick Lewis, Ethan Perez, et al.
\newblock Retrieval-augmented generation for knowledge-intensive NLP tasks.
\newblock \emph{NeurIPS}, 2020.

\end{thebibliography}

\end{document}
